\documentclass[11pt]{article}
\usepackage{amsmath,amssymb,amsthm,hyperref}
\newtheorem{lemma}{Lemma}
\newtheorem{proposition}[lemma]{Proposition}
\newcommand{\F}{\mathbb F}
\newcommand{\gen}[1]{\langle #1\rangle}
\title{Finite end-subgroup certificates for eleven marked tuples}
\author{Portable certificate presentation; attributed classification input}
\date{November 2026}
\begin{document}\maketitle
\paragraph{Scope and attribution.}
Let $G$ be the named inner group $G_2(11)$, acting faithfully on the
seven-dimensional trace-zero split Cayley module over $\F_{11}$.
The identification, common frame and memberships of the four supplied
marks are imported premises. They are not recognized by the checker.
The complete maximal-subgroup cover below, including its natural-module
geometry, is the borrowed input of the retained common note
\texttt{g2\_11\_generation\_by\_prime\_witnesses.tex}, SHA256
\texttt{f1b6c8a69780c7b8ac88f239627a7ce6889c24203f72ddc35c86fe93e661e76b},
and its distinct review
\texttt{8de6c659c024cd16d90bde4f8a8d548735b201b5a5892c747f63d9a8e08ae053}.
The present arguments are classical consequences, not new classification.

\paragraph{Complete-cover premise.}
For $q=11$, the maximal types are the two parabolics,
$(\mathrm{SL}_2(11)\circ\mathrm{SL}_2(11)).2$,
$\mathrm{SL}_3(11):2$, $\mathrm{SU}_3(11):2$,
$2^3.\mathrm{PSL}_3(2)$, $\mathrm{PGL}_2(11)$, $G_2(2)$ and $J_1$.
This is the complete classification list before character-degree
exclusions in Nguyen~\cite[proof of Theorem4.4, pp.10--11]{N}, also
recorded in Wilson~\cite[Table4.1]{W}, attributed to Kleidman~\cite{K}.
There is no proper subfield at the prime field11. The characteristic2
and3 exceptions do not apply. The exceptional prime-field restrictions
for $\mathrm{PSL}_2(13)$ and $\mathrm{PSL}_2(8)$ are checked in
Burness--Thomas~\cite[Proposition7.4]{BT}; even their larger envelopes
introduce neither19 nor37. The original Kleidman publisher full text
was not obtained for the retained assessment; the identified complete
list and independent retained Wilson text are the explicit sources.

The primes of the parabolic and involution-centralizer orders lie in
$\{2,3,5,11\}$. The linear type has primes
$\{2,3,5,7,11,19\}$ and the unitary type has
$\{2,3,5,11,37\}$. The remaining types have respectively prime sets
$\{2,3,7\}$, $\{2,3,5,11\}$, $\{2,3,7\}$ and
$\{2,3,5,7,11,19\}$. Thus no proper maximal has both19 and37,
and only the unitary type has37. In the natural Cayley module that
unitary normalizer is the stabilizer of a nonisotropic rational line
\cite[Section4.2.6 in the retained preliminary edition]{W}.

\begin{proposition}[Two prime words]
If $H\leq G$ contains $X\ne I$, $X^{19}=I$ and $Y\ne I$, $Y^{37}=I$,
then $H=G$.
\end{proposition}
\begin{proof}
Primality gives exact orders19 and37. If $H<G$, a maximal overgroup
would have both primes in its order by Lagrange's theorem, contradicting
the complete-cover premise. No irreducibility or IP hypothesis is used.
\end{proof}

\begin{lemma}[Unique rational line]
If $Y\in\mathrm{GL}_7(\F_{11})$ has order37, its fixed space has
dimension1 and is its unique invariant rational line.
\end{lemma}
\begin{proof}
The order of11 modulo37 is6: $11^2\equiv10$ and
$11^3\equiv-1\pmod{37}$, while $11\not\equiv1$.
Since37 is coprime to11, $x^{37}-1$ is separable. Every nontrivial
irreducible constituent has degree6. A nonidentity order37 operator
in dimension7 therefore has one such degree6 constituent and one
trivial degree1 constituent. A rational invariant line has scalar
action of order dividing both37 and $|\F_{11}^{\times}|=10$, hence is
fixed. The fixed space is exactly the stated1-space.
\end{proof}

\begin{proposition}[One37 word and a moved line]
Suppose $Y\in H\leq G$ has order37, $v\ne0$, $vY=v$, and some
$h\in H$ has $vh\notin\F_{11}v$. Then $H=G$.
\end{proposition}
\begin{proof}
A proper $H$ lies in a maximal with37 in its order. By the imported
cover and geometry it preserves some rational line. That line is
$Y$-invariant, hence equals $\F_{11}v$ by the lemma. Its movement by
$h$ is a contradiction. The checker proves movement using one nonzero
minor $v_i(vh)_j-v_j(vh)_i$. No norm or eigenspace solver is needed.
\end{proof}

\begin{proposition}[Disjoint-label IP obstruction]
Let $(g_0,g_1,g_2,g_3)$ be the supplied distinct nonidentity involutions,
and let $E$ be labels012 or123. If $\gen{g_i:i\in E}=G$, the marked
tuple fails the intersection property.
\end{proposition}
\begin{proof}
The omitted label $j$ is not in $E$, while $g_j\in G$ by the imported
membership. Thus $\gen{g_j}\cap\gen{g_i:i\in E}=\gen{g_j}\ne1$,
but $\{j\}\cap E=\varnothing$, whose generated subgroup is1.
This is precisely a disjoint-label violation. It is not a claim of
excess in the two consecutive rank3 intersections.
\end{proof}

\paragraph{Finite implementation.}
The checker evaluates each supplied end word from the identity, takes
the explicit cofactor, compares all49 entries to the supplied matrix,
and verifies its prime power and nonidentity. For the line cases it
also checks all seven fixed-vector coordinates and the moved minor.
Four involution equations, three distant commutations and distinct
nonidentity marks are checked in every production case. The words and
line vectors are supplied certificates, not conclusions of a replayed
search. Original row/orbit identification and the abstract/natural-module
imports remain separately attributed. No class count, Aut, duality,
census, whole-job, whole-paper or Lean result follows from this package.

\begin{thebibliography}{9}
\bibitem{K} P.B.Kleidman, The maximal subgroups of the Chevalley groups
$G_2(q)$ with $q$ odd, the Ree groups ${}^2G_2(q)$, and their automorphism
groups, \emph{J.Algebra}117(1988),30--71.
\href{https://doi.org/10.1016/0021-8693(88)90239-6}{doi:10.1016/0021-8693(88)90239-6}.
\bibitem{N} H.N.Nguyen, Irreducible restrictions of Brauer characters
of the Chevalley group $G_2(q)$ to its proper subgroups,
\href{https://arxiv.org/abs/0910.4994}{arXiv:0910.4994v1}.
\bibitem{W} R.A.Wilson, \emph{The Finite Simple Groups}, retained
preliminary text633091 bytes, SHA256
\texttt{d0af6d76130d73a143e76218c7c66ca0167d06f70d591f9c3482ec8e04fe895f},
Sections4.2.5--4.2.6,Table4.1 and Section5.9.1.
The published edition renumbers the Cayley chapter as Section4.3.
\url{https://webspace.maths.qmul.ac.uk/r.a.wilson/TFSG.html}.
\bibitem{BT} T.C.Burness and A.R.Thomas, The classification of extremely
primitive groups, authors' primary manuscript, Proposition7.4.
\url{https://seis.bristol.ac.uk/~tb13602/docs/ep_exceptional_final.pdf}.
\end{thebibliography}\end{document}
